\documentclass{article}
\usepackage{azzam}

\begin{document}

\title{Geometri Dasar}
\maketitle
\section{Latihan Soal Segitiga }
\begin{enumerate}    
    \item Garis berat $AD$ pada segitiga $ABC$ memotong garis berat $CF$ di titik $P$, serta perpanjangan $BP$ memotong $AC$ di $E$. Jika diketahui segitiga $ABC$ lancip dan $AB=6$, maka panjang $DE$ adalah \dots

    \item (OSK 2011,2012,2013,2018) Diberikan segitiga $ABC$ dan lingkaran $\Gamma$ yang berdiameter $AB$. Lingkaran $\Gamma$ memotong sisi $AC$ dan $BC$ berturut-turut di titik $D$ dan $E$. Jika $AD = \frac13 AC, BE =\frac14 BC$ dan $AB = 30$, maka luas segitiga $ABC$ adalah \dots
		
    \item Diberikan segitiga $ABC$ dengan $D$ titik tengah $AC$, $E$ titik tengah $BD$, dan $H$ merupakan pencerminan $A$ terhadap $E$. Jika $F$ merupakan perpotongan antara $AH$ dengan $BC$, maka nilai $\dfrac{AF}{FH}$ sama dengan \dots
		 
    \item Diberikan segitiga $ABC$ dengan panjang sisi $BC = 20$, $CA = 24$, dan $AB=12$. Titik $D$ pada segmen $BC$ dengan $BD = 5$. Lingkaran luar dari segitiga $ABD$ memotong $CA$ di $E$. Hitunglah nilai $2 \times DE$.

    \item (OSK 2015) Diberikan trapesium $ABCD$ dengan $AB$ sejajar $DC$ dan $AB = 84$ serta $DC = 25$. Jika trapesium $ABCD$ memiliki lingkaran dalam yang menyinggung keempat sisinya, keliling trapesium $ABCD$ adalah \ldots

    \item (OSK 2022) Diberikan segitiga siku-siku $ABC$. Jika luas dari segitiga $ABC$ adalah 112. Misalkan $R$ adalah panjang jari-jari lingkaran luar segitiga $ABC$ dan $r$ adalah panjang jari-jari lingkaran dalam segitiga $ABC$. Diketahui juga $R + r = 16$. Panjang sisi miring dari segitiga $ABC$ adalah \ldots
\end{enumerate}
\section{Latihan Soal Pythagoras}
\begin{enumerate}
    \item (OSK SMP 2016) Diketahui $ABCD$ dan $CEGH$ adalah dua persegipanjang kongruen dengan panjang $17$ cm, dan lebar $8$ cm. Titik $E$ berada di sisi $AB$ dan $D$ berada di sisi $GH$. Titik $F$ adalah titik potong sisi $AD$ dan $EG$. Luas segiempat $EFDC$ adalah .... $cm^2$.

    \item Misalkan $ABC$ adalah segitiga lancip. Titik $D$, $E$, dan $F$ terletak pada sisi $BC$, $CA$, dan $AB$, berturut-turut, sedemikian sehingga $AD$, $BE$, dan $CF$ adalah garis tinggi segitiga $ABC$. Titik $H$ adalah titik tinggi segitiga $ABC$. Jika $DE = 8$, $DF = 15$, dan $EF = 17$, tentukan panjang $AH$.

    \item (OSK 2015) Diberikan segitiga $ABC$ dengan sudut $\angle ABC = 90^\circ$. Lingkaran $L_1$ dengan $AB$ sebagai diameter sedangkan lingkaran $L_2$ dengan $BC$ sebagai diameternya. Kedua lingkaran $L_1$ dan $L_2$ berpotongan di $B$ dan $P$. Jika $AB = 5$, $BC = 12$ dan $BP = x$ maka nilai dari $\frac{240}{x}$ adalah \ldots

    \item (OSK 2022) Diberikan segitiga $ABC$ siku-siku di B. Titik $D$ berada pada sisi $AB$ dan titik $E$ berada pada sisi $AC$. Diketahui $DE$ sejajar $BC$. Jika $AD = 21$, $DB = 3$, dan $BC = 32$, maka panjang $AE$ adalah \dots

\end{enumerate}
\section{Latihan Soal Lingkaran}
\begin{enumerate}
    \item Pada segiempat $WXYZ$ dengan diagonal yang saling tegak lurus diketahui bahwa $\angle WZX = 30^\circ, \angle XWY = 40^\circ,$ and $\angle WYZ = 50^\circ$. Hitunglah besar $\angle X$ dan $\angle Z$.
    \begin{center}
        \includegraphics[width=0.25\textwidth]{0Figure/evanquad.PNG}
    \end{center}
		
    \item (OSK 2013) Diberikan segitiga lancip $ABC$ dengan $O$ sebagai pusat lingkaran luarnya. Misalkan $M$ dan $N$ berturut - turut pertengahan $OA$ dan $BC$. Jika $\angle ABC = 4\angle OMN$ dan $\angle ACB = 6\angle OMN$, maka besarnya $\angle OMN$ sama dengan \dots

    \item 	Pada gambar di bawah, diketahui titik A $\ne$ B pada lingkaran berdiameter $MN$ dan berpusat di $C$. $P$ adalah titik pada segmen $CN$ dimana $\angle CAP = \angle CBP = 10 ^\circ$. Jika $\angle ACM = 40^\circ$, maka $\angle BCN = \dots^\circ$	
    \begin{center}
         \includegraphics[width=0.25\textwidth]{0Figure/soalLingkaran1.PNG}
    \end{center}

    \item (OSK 2011,2012,2013,2018) Diberikan segitiga $ABC$ dan lingkaran $\Gamma$ yang berdiameter $AB$. Lingkaran $\Gamma$ memotong sisi $AC$ dan $BC$ berturut-turut di titik $D$ dan $E$. Jika $AD = \frac13 AC, BE =\frac14 BC$ dan $AB = 30$, maka luas segitiga $ABC$ adalah \dots

    \item (OSK 2015) Diberikan segitiga $ABC$ dengan sudut $\angle ABC = 90^\circ$. Lingkaran $L_1$ dengan $AB$ sebagai diameter sedangkan lingkaran $L_2$ dengan $BC$ sebagai diameternya. Kedua lingkaran $L_1$ dan $L_2$ berpotongan di $B$ dan $P$. Jika $AB = 5$, $BC = 12$ dan $BP = x$ maka nilai dari $\frac{240}{x}$ adalah \ldots
\end{enumerate}
\section{Latihan Soal Luas}
\begin{enumerate}
\item (OSK 2006) Pada segitiga $ABC$, titik $F$ membagi sisi $AC$ dalam perbandingan $1 : 2$. Misalkan $G$ titik tengah $BF$ dan $E$ titik perpotongan antara sisi $BC$ dengan $AG$. Maka titik $E$ membagi sisi $BC$ dalam perbandingan \dots

\item (OSK 2009) Diberikan persegi $ABCD$ dengan panjang sisi 10. Misalkan $E$ pada $AB$ dan $F$ pada $BC$ dengan $AE = FB = 5$. Misalkan $P$ adalah titik potong $DE$ dan $AF$. Luas $DCFP$ adalah \dots

\item Pada segitiga $ABC$, titik $D$ dan $E$ berturut-turut berada di segmen $AB$ dan $AC$. Misalkan $CD$ dan $BE$ berpotongan di titik $F$. Jika $[DBF]=3$, $[BFC]=6$, dan $[CFE]=4$, berapakah luas $ADFE$?

\item (OSK 2016) Pada segitiga ABC, titik $X, Y$ dan $Z$ berturut-turut terletak pada sinar $BA, CB$ dan $AC$ sehingga $BX = 2BA, CY = 2CB$ dan $AZ = 2AC$. Jika luas segitiga $ABC$ adalah 1, maka luas segitiga $XYZ$ adalah \dots

\item (OSK 2016) Segitiga $ABC$ merupakan segitiga sama kaki dengan panjang $AB = AC = 10 $. Titik $D$ terletak pada garis $AB$ sejauh $7 $ dari $A$ dan $E$ titik pada garis $AC$ yang terletak sejauh $4 $ dari $A$. Dari $A$ ditarik garis tinggi dan memotong $BC$ di $F$. Jika bilangan rasional $\frac{a}{b}$ menyatakan perbandingan luas segi empat $ADFE$ terhadap luas segitiga $ABC$ dalam bentuk yang paling sederhana, maka nilai $a + b$ adalah \dots

\item (OSK 2014) Diberikan segitiga $ABC$ yang sisi-sisinya tidak sama panjang sehingga panjang garis berat $AN$ dan $BP$ berturut-turut 3 dan 6. Jika luas segitiga $ABC$ adalah $3\sqrt{15}$, maka panjang garis berat ketiga $CM$ adalah \dots

\item (OSK 2012) Diketahui $\triangle ABC$ sama kaki dengan panjang $AB = AC = 3$, $BC = 2$, titik $D$ pada sisi $AC$ dengan panjang $AD = 1$. Tentukan luas $\triangle ABD$.

\item (OSK 2012) Diberikan segitiga $ABC$ dengan keliling 3, dan jumlah kuadrat sisi-sisinya sama dengan 5. Jika jari-jari lingkaran luarnya sama dengan 1, maka jumlah ketiga garis tinggi dari segitiga $ABC$ tersebut adalah \dots

\item (OSN 2011 SMP/MTs) Bangun datar $ABCD$ adalah trapesium dengan $AB$ sejajar $CD$ dan $AB < CD$. Titik $E$ dan $F$ terletak ada $CD$ sehingga $AD$ sejajar $BE$ dan $AF$ sejajar $BC$. Titik $H$  adalah perpotongan $AF$ dengan $BE$ dan titik $G$ adalah  perpotongan $AC$ dengan $BE$. Jika panjang $AB$ adalah 4 cm dan panjang $CD$ adalah 10 cm hitunglah perbandingan luas  segitiga $AGH$ dengan luas trapesium $ABCD$. 
\end{enumerate}
\section{Latihan Soal Power of A Point}
\begin{enumerate}
    \item (\textbf{Soal Legend: OSK 2011,2012,2013,2018}) Diberikan segitiga $ABC$ dan lingkaran $\Gamma$ yang berdiameter $AB$ . Lingkaran $\Gamma$ memotong sisi $AC$ dan $BC$ berturut-turut di titik $D$ dan $E$. Jika $AD = \frac13 AC, BE =\frac14 BC$ dan $AB = 30$, maka luas segitiga $ABC$ adalah \dots

    \item (OSN SL 2010) Pada segitiga $ABC$, misalkan $D$ adalah titik tengah $BC$, dan $BE$, $CF$ adalah garis tinggi. Buktikan bahwa $DE$ dan $DF$ keduanya adalah garis singgung lingkaran luar $\triangle AEF$. %(OSN SL 2010) 

    %AIME 2016 I and AIME 2016 I
    \item (AIME 2016 I) Misalkan $\triangle ABC$ adalah segitiga lancip dengan lingkaran $\omega,$ dan misalkan $H$ adalah titik potong dari garis tinggi $\triangle ABC.$ Garis singgung lingkaran luar $\triangle HBC$ di $H$ memotong $\omega$ pada titik $X$ dan $Y$ dengan $HA=3,HX=2,$ dan $HY=6.$ Carilah luas dari $\triangle ABC$.

    \item (AIME 2016 I) Lingkaran $\omega_1$ dan $\omega_2$ bertemu di titik $X$ dan $Y$. Garis $\ell$ menyinggung lingkaran $\omega_1$ dan $\omega_2$ di $A$ dan $B$, berturut-turut, dengan garis $AB$ lebih dekat ke titik $X$ daripada $Y$. Lingkaran $\omega$ yang melewati $A$ dan $B$, memotong $\omega_1$ lagi di $D \neq A$ dan memotong $\omega_2$ lagi di $C \neq B$. Ketiga titik $C$, $Y$, $D$ segaris dengan $XC = 67$, $XY = 47$, dan $XD = 37$. Carilah panjang $AB$.
\end{enumerate}
\section{Latihan Soal Dalil Ceva}
\begin{enumerate}
    \item (OSK 2015) Pada segitiga $ABC$, garis tinggi $AD$, garis bagi $BE$ dan garis berat $CF$ berpotongan di satu titik. Jika panjang $AB = 4$ dan $BC = 5$, dan $CD = \frac{m^2}{n^2}$ dengan $m$ dan $n$ relatif prima, maka nilai dari $m - n$ adalah \ldots
\end{enumerate}
\section{Latihan Soal Dalil Menelaus}
\begin{enumerate}
    \item Dalam segitiga $ABD$, $F$ berada pada segmen $AD$, $E$ berada pada sinar $BF$, $G$ berada pada segmen $BD$, dan $C$ adalah titik perpotongan dari $FG$ dan $ED$. Diketahui bahwa $AB = 15$, $BD = 18$, $AF = 15$, $DF = 12$, $BE = 24$, dan $CF = 17$. Temukan rasio $BG : FG$.
    % https://drive.google.com/drive/search?q=parent:0B-4OltLGFEDFeG92VENBbmdZODA%20type:pdf%20menelaus

    \item (OSK 2022) Diberikan $ABC$ siku-siku sama kaki dengan $BC=AB$. Misalkan $L$ titik tengah $BC$ dan $P$ pada sisi $AC$ sehingga $BP \perp AL$. Jika $CP=30\sqrt{2}$, maka panjang $AB$ adalah \ldots

    \item Diberikan segitiga $ABC$ dengan panjang $BC = 36$. Misalkan $D$ adalah titik tengah $BC$ dan $E$ adalah titik tengah $AD$. Misalkan pula bahwa $F$ adalah perpotongan $BE$ dengan $AC$. Jika diketahui bahwa $AB$ menyinggung lingkaran luar segitiga $BFC$, hitunglah panjang $BF$.
\end{enumerate}
\section{Latihan Soal Dalil Ptolemy}
\begin{enumerate}
    \item Diberikan sebuah segiempat siklis $ABCD$ dengan $ABC$ adalah segitiga sama sisi. Jika $AD=2$ dan $CD=3$, panjang $BD=\dots$
\end{enumerate}
\section{Latihan Soal Dalil Stewart}
\begin{enumerate}
    \item (OSK 2016) Pada segitiga $ABC$, titik $M$ terletak pada $BC$ sehingga $AB=7, AM=3, BM=5$, dan $MC=6$. Panjang $AC$ adalah \dots

    \item (HMMT 1999) Dalam segitiga $ABD$, $F$ berada pada segmen $AD$, $E$ berada pada sinar $BF$, $G$ berada pada segmen $BD$, dan $C$ adalah titik perpotongan dari $FG$ dan $ED$. Diketahui bahwa $AB = 15$, $BD = 18$, $AF = 15$, $DF = 12$, $BE = 24$, dan $CF = 17$. Temukan rasio $BG : FG$.
    % https://drive.google.com/drive/search?q=parent:0B-4OltLGFEDFeG92VENBbmdZODA%20type:pdf%20menelaus
\end{enumerate}
\section{Latihan Soal Teorema Garis Bagi}
\begin{enumerate}
    \item (OSK 2014) Diberikan segitiga $ABC$ dengan $AB = 360, BC = 240,$ dan $AC = 180$. Garis bagi dalam dan garis bagi luar dari $\angle CAB$ memotong $BC$ dan perpanjangan $BC$ berturut-turut di $P$ dan $Q$. Jari-jari lingkaran yang melalui titik-titik $A, P,$ dan $Q$ adalah \dots
    
    \item (OSK 2015) Pada segitiga $ABC$, garis tinggi $AD$, garis bagi $BE$ dan garis berat $CF$ berpotongan di satu titik. Jika panjang $AB = 4$ dan $BC = 5$, dan $CD = \frac{m^2}{n^2}$ dengan $m$ dan $n$ relatif prima, maka nilai dari $m - n$ adalah \ldots
\end{enumerate}
\section{Latihan Soal Miquel Theorem}
\begin{enumerate}
    \item (OSK 2016) Diberikan empat titik pada satu lingkaran $\Gamma$ dalam urutan $A,B,C,D$. Sinar garis $AB$ dan $DC$ berpotongan di $E$, dan sinar garis $AD$ dan $BC$ berpotongan di $F$. Misalkan $EP$ dan $FQ$ menyinggung lingkaran $\Gamma$ berturut-turut di $P$ dan $Q$. Misalkan pula bahwa $EP=60$ dan $FQ=63$, maka panjang $EF$ adalah \dots
\end{enumerate}
\section{Latihan Soal Trigonometri}
\begin{enumerate}
    \item (OSK 2022) Diberikan segitiga $ABC$ seperti di gambar, dengan panjang $AB = 2BC$ dan $BD = CD$. Jika luas segitiga $DEC$ adalah 10, luas dari segitiga $AFE$ adalah \dots

	\item Jika $A+B=45^\circ$ dan $\cos A\sin B=\frac{\sqrt{2}}{6}$, maka $\cos(B-A)=\dots$
	
	\item Nilai dari $\cos \dfrac{\pi}{7}\cdot \cos \dfrac{2\pi}{7} \cdot \cos \dfrac{4\pi}{7}$ adalah \dots
	
	\item Pada segitiga $ABC$, buktikan bahwa $\tan A + \tan B + \tan C = \tan A \tan B \tan C$.
	
	\item Tentukan nilai eksak dari $\tan 1^\circ \cdot \tan 2^\circ \cdot \tan 3^\circ \cdot \ldots \cdot \tan 89^\circ$.
	
	\item (OSK 2005) Nilai dari $\sin^8 75^\circ - \cos^8 75^\circ$ adalah \dots
\end{enumerate}
\section{Latihan Soal Dalil Sinus dan Cosinus}
\begin{enumerate}
    \item (OSK 2016) Pada segitiga $ABC$, titik $M$ terletak pada $BC$ sehingga $AB=7, AM=3, BM=5$, dan $MC=6$. Panjang $AC$ adalah \dots

    \item (OSK 2013) Misalkan $P$ adalah titik interior dalam daerah segitiga $ABC$ sehingga besar $\angle PAB = 10^\circ, \angle PBA = 20^\circ, \angle PCA = 30^\circ, \angle PAC=40^\circ$. Besar $\angle ABC = \dots$
    
    \item (LMNAS SMP 33 Penyisihan) Diberikan segitiga tumpul $ABC$ dengan $AB = BC$. Titik $D$ berada di dalam segitiga tersebut sedemikian sehingga $AD = BD$, $\angle ADB = 140^\circ$, dan $\angle ADC = 150^\circ$. Besar sudut $ACD$ dalam satuan $^\circ$ (derajat) adalah \dots
    
    \item (Modifikasi OSK 2017) Pada sebuah lingkaran dengan pusat $O$, talibusur $AB$ berjarak 5 dari titik $O$ dan talibusur $AC$ berjarak $5\sqrt{2}$ dari titik $O$ dengan titik $A$ terletak di busur $BC$ yang lebih kecil ($A$ diantara $B$ dan $C$) Jika panjang jari-jari lingkaran 10, maka $BC^2=\dots$
\end{enumerate}
\end{document}